\origpage{252}{ROBERTSON AND SEYMOUR}
\noindent are disjoint connected subgraphs $H_1, \dots, H_p$ of $\Gamma$ with $V(H_j) \cap Z = Z_j$ ($1 \le j \le p$).

For $1 \le i < j \le p$, let $Q_{ij}$ be the path obtained from $P_{ij}$ by deleting $v_{ij}$, and let $e_{ij}$ be the edge of $P_{ij}$ incident with $v_{ij}$. For $1 \le j \le p$, let $X_j = H_j \cup \bigcup (Q_{ij} \colon 1 \le i < j)$. We see that $X_1, \dots, X_p$ are mutually disjoint connected subgraphs of $G$, and for $1 \le i < j \le p$ there is an edge of $G$ between $V(X_i)$ and $V(X_j)$, namely, $e_{ij}$.

Let $\mathscr{T}^*$ be the tangle in $G$ induced by $\mathscr{T}$, let $(A, B) \in \mathscr{T}^*$, and suppose that $V(X_j) \subseteq V(A)$ for some $j$. Since $p \ge 5 \ge 3$, it follows that $|Z_j| \ge 2$; choose $u, v \in Z_j$, distinct. Then $d(u, v) \ge \theta$. Now $u, v \in Z_j \subseteq V(H_j)$, and $H_j$ is connected; let $P$ be a path of $H_j$ with ends $u, v$. Then $V(P) \subseteq V(H_j) \subseteq V(A \cap \Gamma)$, and $(A \cap \Gamma, B \cap \Gamma) \in \mathscr{T}$ by definition of $\mathscr{T}^*$, and so by (4.1), $(A \cap \Gamma, B \cap \Gamma)$ has order $\ge \frac{1}{2} d(u, v) \ge \frac{1}{2} \theta \ge p$, and hence so does $(A, B)$. It follows that $\mathscr{T}^*$ controls a $K_p$ minor. $\square$

\stmt
(4.5) \textit{Let $\Sigma$ be a connected surface, and let $p \ge 0$ be an integer. Then there are integers $\theta \ge 1$ and $t \ge 0$ such that the following is true. Let $\Gamma$ be a 2-cell drawing in $\Sigma$, with a respectful tangle $\mathscr{T}$ of order $\ge \theta$, with metric $d$. Let $r_1, \dots, r_t$ be distinct regions of $\Gamma$ such that $d(r_i, r_j) \ge \theta$ for $1 \le i < j \le t$. For $1 \le i \le t$, let $C_i$ be a circuit of $\Gamma$ such that $U(C_i)$ bounds $r_i$; and let $a_i, b_i, c_i, d_i \in V(C_i)$ be distinct, so that $a_i, b_i, c_i, d_i$ occur in $V(C_i)$ in that order, and $\{a_i, b_i, c_i, d_i\}$ is free with respect to $\mathscr{T}$. Let $\Gamma$ be a subgraph of a graph $G$, and for $1 \le i \le t$ let $P_i, Q_i$ be paths of $G$ with ends $a_i, c_i$ and $b_i, d_i$, respectively, and with no other vertices or edges in $\Gamma$. Let $P_1, Q_1, \dots, P_t, Q_t$ be mutually disjoint. Then the tangle in $G$ induced by $\mathscr{T}$ controls a $K_p$ minor.}

\textit{Proof.} As in (4.4), we may assume that $p \ge 5$. Now $K_p$ can be drawn with crossings in the plane, and hence in $\Sigma$, so that (i) no edge passes through any vertex, (ii) every two edges meet only finitely many times, (iii) whenever two edges meet they cross (that is, edges never touch ``tangentially''), (iv) no three edges pass through any point of the surface, (v) no edge crosses itself, (vi) every edge crosses at least one other edge, and (vii) every vertex is incident with two edges that do not cross each other. To see this, first draw a Hamilton circuit without crossings, and then add the remaining edges to the drawing, so that (i)--(vi) are satisfied. Then (vii) is satisfied too.

Let this drawing have $t$ crossings. Let $L$ be obtained from $K_p$ by subdividing its edges, where if $e \in E(K_p)$ is involved in $k$ crossings of the drawing of $K_p$ then $e$ is subdivided $k - 1$ times. It follows that $L$ can be drawn in $\Sigma$ so that every edge is involved in exactly one crossing, and (i)--(vii) above are satisfied. Thus $t = \frac{1}{2} |E(L)|$, and the edges of $L$ fall into $t$ crossing pairs. Let the crossings be $x_1, \dots, x_t$. By choosing a small disc around each crossing and applying a suitable homeomorphism, we deduce
